% =====================================================================
% Conclusion
% Status: careful placeholder. No empirical claim is written until
% the pipeline has produced results.
% =====================================================================
\section{Conclusion}\label{sec:conclusion}

This thesis has asked whether there is seasonal variation in the
spatial extremal dependence of extreme wind gusts at KNMI station
locations in the Netherlands. The empirical strategy has been to
estimate, separately for winter (DJF) and summer (JJA), the
pairwise upper-tail dependence coefficient
\(\chi_{ij}^{(c)}\) and the extremal coefficient
\(\theta_{ij}^{(c)}\) for every station pair on a balanced subset
of the KNMI network, and to summarise these along the geographical
distance \(d_{ij}\) between stations. The seasonal comparison is
then a comparison of the two dependence--distance curves
\(\widehat\chi^{(\mathrm{W})}(d)\) and
\(\widehat\chi^{(\mathrm{S})}(d)\) on the same panel.

% TODO: this section currently states the question and the
%       methodology only. The numerical findings will be inserted
%       after the empirical pipeline of Section~\ref{sec:methodology}
%       has been run on the actual data. Do not pre-write empirical
%       conclusions.

\subsection{Summary of findings}\label{sec:conclusion-summary}

% TODO: insert 1--2 short paragraphs that answer the research
%       question, using careful language ("evidence suggests"
%       rather than "proves"). Refer to the dependence--distance
%       curves and the scalar summaries at the reference
%       distances. Do not draw climate-change conclusions.

\subsection{Limitations}\label{sec:conclusion-limitations}

The thesis has three first-order limitations.

\begin{enumerate}
    \item The changing KNMI station network means that any
          temporal comparison (the secondary sample-split analysis
          in Section~\ref{sec:rob-period}) cannot cleanly separate
          changes in the dependence structure from changes in
          spatial coverage. The seasonal comparison is robust to
          this concern because winter and summer are estimated on
          the same balanced panel.
    \item The finite-sample tail estimators of
          Section~\ref{sec:meth-dependence} report on
          \(\chi_{ij}^{(c)}(u)\) at finite \(u\), not at the
          asymptotic limit. \citet{TODO-wadsworth2012} show that
          finite-level dependence can decay to zero at the very
          top of the tail; the thesis addresses this only
          partially by reporting on a grid of \(u\) and by
          including the \(\bar\chi\) diagnostic.
    \item The dependence--distance summary assumes isotropy of the
          dependence. Anisotropy due to prevailing wind direction
          is a plausible second-order effect that the thesis
          examines only as a robustness check
          (Section~\ref{sec:rob-direction}).
\end{enumerate}

\subsection{Directions for further work}\label{sec:conclusion-future}

Three directions are natural extensions of the work undertaken
here.

\begin{enumerate}
    \item Fit a covariate-dependent (season-indexed) max-stable
          process in the spirit of
          \citet{TODO-huser-genton2016}, using pairwise
          composite likelihood. This would replace the
          ``estimate separately by season'' strategy of the
          thesis with a single non-stationary model and would let
          the seasonal contrast be tested directly through a
          likelihood-ratio statistic.
    \item Reformulate the analysis at the level of intraday gust
          curves rather than hourly maxima, using the functional
          peaks-over-threshold framework of
          \citet{TODO-fondeville2022} or the spatio-temporal
          functional construction of \citet{TODO-french2018}.
          This would allow the spatial dependence of \emph{whole
          storm footprints} rather than of single hours to be
          estimated.
    \item Extend the analysis from pairwise dependence to joint
          high-dimensional dependence using vine copulas in the
          construction of \citet{TODO-aas2008}, treating the
          present pairwise summaries as the building blocks. Vine
          copulas were excluded as the main method of this thesis
          and are recorded here only as a future-work direction.
\end{enumerate}

A more modest extension is the verification simulation described
in Section~\ref{sec:simulation}: it is optional and not part of
the main contribution, but it is the cleanest available check
that the estimator behaves as expected under known dependence
regimes.
